\documentclass[11pt]{article}
\usepackage{amsmath,amssymb}
\usepackage[margin=1in]{geometry}
\usepackage{booktabs}
\usepackage{tabularx}
\usepackage{hyperref}
\setlength{\emergencystretch}{2em}

\title{General Scientific Claim-Authorization Superiority:\\
When Correct Evidence Actually Licenses a Scientific Conclusion}
\author{ORION Paper IV Ultimate-Successor Working Manuscript}
\date{August 2026}

\begin{document}
\maketitle

\begin{abstract}
Modern agent systems already support provenance tracing, partial-evidence awareness, abstention, evidence-backed permission graphs, formal authorization and execution-policy verification. ORION Paper IV therefore targets a higher scientific responsibility: across heterogeneous domains, evidence structures, state representations and adversarial benchmark constructions, can ORION outperform a donor-complete system at deciding when evidence actually licenses a scientific conclusion, while preserving valid conclusions and avoiding excessive rejection under matched resources? The predecessor H1/H2 evidence remains immutable, and the saturated H3 slice remains \textsc{NotSupported} because its construction did not identify the intended competence. P4-U turns that null into a benchmark-science and mechanism-discovery programme. Strong comparators receive modern provenance, verification, partial-evidence reporting, abstention, AuthGraph/FAVA-style authorization and responsibility-aware state. The residual is scientific promotion authority over evidence sufficiency, source dependence, statistical validity, scope, time, reproducibility, custody and evaluator validity. The intended terminal is \textbf{GENERAL\_SCIENTIFIC\_AUTHORIZATION\_SUPERIORITY}, including at least one failure-derived promotion obligation or benchmark-identifiability mechanic that transfers to fresh constructions.
\end{abstract}

\section{Largest scientific question}
\begin{quote}
Across heterogeneous scientific evidence situations, can ORION decide more accurately than the strongest donor-complete system whether a candidate scientific conclusion is licensed, and can it preserve that advantage after benchmark shortcuts, representation confounds and evaluator attacks are removed?
\end{quote}

Partial Evidence Bench directly measures unsafe completeness when material evidence lies outside an authorized view \cite{partial2026}. AuthGraph structurally aligns execution provenance with an isolated authorization graph \cite{authgraph2026}. FAVA lowers natural-language constraints into evidence-backed permission graphs and uses deterministic authorization before effectful action \cite{fava2026}. CVA formalizes cryptographically verifiable agent authorization binding principal, request, context and policy \cite{cva2026}. These systems raise the donor baseline. P4-U does not claim provenance or authorization itself as novel. It asks whether ORION can generalize over these mechanisms to the distinct scientific question of conclusion promotion.

\section{Upward-generalization principle}
Donor pressure expands the comparator rather than shrinking the final target. ORION earns the broad claim only through a combination of lower false promotion, preserved clean promotion, correct responsibility-rung use, construction-identifiability robustness, cross-domain transfer and failure-driven expansion of the promotion-obligation basis.

The intended terminal is \textbf{GENERAL\_SCIENTIFIC\_AUTHORIZATION\_SUPERIORITY}.

\section{Immutable predecessor evidence}
The predecessor P4 protected V2 results are unchanged. H1 remains the registered false-promotion result, H2 the clean-coverage guard, and H3 remains \textsc{NotSupported} because the frozen H3 construction was non-identifying. No old label, denominator, margin or interpretation is rewritten.

The scientific lesson is stronger than a local null: a benchmark score cannot authorize a scientific competence claim until the benchmark itself is shown to identify that competence rather than a construction shortcut.

\section{Scientific authority as responsibility}
For evidence state $E$, claim $C$, context $X$ and requested responsibility $R$,
\[
A(E,C,X,R)\in\{\textsc{Authorized},\textsc{Blocked},\textsc{Unresolved},\textsc{CannotCheck}\},
\]
with at least
\[
R\in\{\textsc{Answer},\textsc{Verify},\textsc{Promote},\textsc{Repair}\}.
\]
A state sufficient for \textsc{Answer} may be insufficient for \textsc{Promote}. Correct final text is therefore not evidence of promotion competence.

\section{Construction-balanced benchmark}
Before protected outcomes, construction, labels, margins and shortcut thresholds are frozen. Labels must not be recoverable from source count, evidence-list length, missingness, template, field order, path name or metadata. Mandatory families include partial evidence, conflicting evidence, hidden source dependence, invalid statistics, unreplayable code, contamination, scope mismatch, temporal invalidation, evaluator defects, responsibility mismatch and clean fully supported positives.

At least two independently generated benchmark constructions and multiple scientific domains are required for the headline claim.

\section{Donor-complete baselines}
All systems receive matched model, tools, evidence access, hidden-field restrictions and total resource budgets. Required comparators include:
\begin{enumerate}
    \item factual verification and calibrated abstention;
    \item provenance/attribution and claim-evidence graph systems;
    \item Partial-Evidence-Bench-style completeness/gap reporting \cite{partial2026};
    \item AuthGraph-style provenance/authorization alignment \cite{authgraph2026};
    \item FAVA-style evidence-backed permission graphs and formal authorization \cite{fava2026};
    \item cryptographically verifiable authorization as an analysis-strength donor \cite{cva2026};
    \item contamination and benchmark-audit systems;
    \item responsibility-aware state/certificate systems;
    \item donor-complete product containing all preceding mechanisms and allowed to route among them;
    \item strongest runnable scientific verification/claim-checking system at protocol freeze;
    \item oracle authority ceiling for analysis only.
\end{enumerate}

\section{Scientific Promotion Obligation Discovery}
The first ChatGPT-led negative-to-positive successor family is \textbf{Scientific Promotion Obligation Discovery (SPOD)}. Every retained false promotion, false rejection, null or benchmark failure is treated as competing responsibility hypotheses over evidence sufficiency, provenance, dependence, statistics, scope, temporal validity, reproducibility, state responsibility, contamination/custody, evaluator validity, benchmark identifiability and the authority rule itself.

SPOD freezes one discriminator that changes a suspected obligation while preserving answer content and other coordinates. A new promotion obligation is admitted only when strong verifier/provenance/authorization donors fail prospectively, an independently grounded evaluator identifies the missing obligation, and the resulting rule transfers to a fresh domain/construction without harming clean promotion. Generic policy learning, provenance graphs and authorization graphs remain donor-owned.

\section{Primary hypotheses}
\paragraph{H1, authorization superiority.} ORION reduces false scientific promotion versus the strongest runnable donor-complete comparator by a frozen margin.

\paragraph{H2, coverage preservation.} H1 is invalid if valid clean promotion falls below the frozen non-inferiority guard or false rejection rises beyond its guard.

\paragraph{H3, identifiability robustness.} The advantage survives length/source-count/missingness/template/metadata attacks and independently generated balanced reconstructions.

\paragraph{H4, responsibility-rung superiority.} ORION more accurately distinguishes answer, verify, promote and repair sufficiency.

\paragraph{H5, mechanism expansion.} At least one retained null/failure yields a new promotion obligation or benchmark-identifiability mechanic that transfers beyond its originating construction.

\section{Statistics and hostile controls}
The protected scientific case or prospectively declared construction cluster is the independent unit. Repeated generations and remints are technical repeats. Freeze superiority/non-inferiority margins, cluster structure, multiplicity, interval method, catastrophic false-promotion rule and \textsc{CannotCheck} handling before protected access.

False promotion, false rejection, clean coverage, calibration, responsibility-rung errors and shortcut susceptibility are reported separately. Required hostile controls include correct-answer/unauthorized-evidence pairs, wrong-scope and stale evidence, hidden dependence, self-authored/stale evaluators, contamination, answer-sufficient but promotion-insufficient state, always-abstain, always-open-raw, and donor-complete routing.

\section{Framework correspondence and external authority}
The executable framework supplies a strong evidence-authority substrate through \texttt{ExternalEvidenceManifest.v1}. Each external record binds the exact subject revision, evidence artifact, evaluator artifact, producer and verifier process lineages, evaluation epoch, split, frozen-before-candidate chronology and freshness. Missing, duplicate, mismatched, self-verified, post-hoc or non-fresh evidence yields \textsc{CannotCheck}; a verified \textsc{Fail} remains \textsc{Fail}.

P4-specific external criteria already require source-attribution benchmarking, contamination audit, evaluator locking, held-out access logging, a matched verifier baseline and false-promotion improvement. These mechanics make P4-U semantically consistent with the code base but do not themselves establish general scientific authorization superiority.

SPOD, naturalistic scientific promotion inference, automatic benchmark repair and the headline terminal remain prospective and are not presented as registered runtime capability.

\section{Claim ladder and current status}
The progression is: protocol frozen; protected authorization superiority; identifiability-robust superiority; cross-domain/construction replication; and \textbf{GENERAL\_SCIENTIFIC\_AUTHORIZATION\_SUPERIORITY} with H5 mechanism expansion. This TeX is prospective. It preserves the predecessor H3 null and asserts no general superiority until external evidence satisfies the canonical manifest.

\begin{thebibliography}{8}
\bibitem{partial2026}
K. Tallam, ``Partial Evidence Bench: Benchmarking Authorization-Limited Evidence in Agentic Systems,'' arXiv:2605.05379, 2026.

\bibitem{authgraph2026}
P. Wang, Y. Li, and Y. Tian, ``Aligning Provenance with Authorization: A Dual-Graph Defense for LLM Agents,'' arXiv:2605.26497, 2026.

\bibitem{fava2026}
Y. Zhang, X. Zhao, S. Liu, H. Lou, G. Cheng, and C. Liu, ``FAVA: Formal Authorization for Verified Agents with Evidence-Backed Permission Graphs,'' arXiv:2607.27267, 2026.

\bibitem{cva2026}
M. Llamb\'i-Morillas and D. Fern\'andez-Fern\'andez, ``Toward cryptographically verifiable authorization for autonomous AI agents: A security hypothesis, preliminary formal model, and proof-of-concept implementation,'' arXiv:2607.21325, 2026.
\end{thebibliography}

\end{document}
